\documentclass[11pt]{article}

\usepackage[margin=1in]{geometry}
\usepackage[T1]{fontenc}
\usepackage{lmodern}

\usepackage{amsmath, amssymb}

\usepackage{setspace}

\usepackage[most]{tcolorbox}
\usepackage{xcolor}

\tcbset{
  enhanced,
  boxrule=0pt,
  frame hidden,
  breakable
}

\usepackage{fancyhdr}
\pagestyle{fancy}
\fancyhf{}
\setlength{\headheight}{52pt}
\renewcommand{\headrulewidth}{0pt}

\newcommand{\Course}{Math 4610 Spring 2026}
\newcommand{\Instructor}{Discussion 2}
\newcommand{\AssignmentType}{Discussion} % or Discussion
\newcommand{\AssignmentNumber}{2}

\newcommand{\Contributors}{
Marks, Stephen\\
}


\fancyhead[L]{\Course\\ \Instructor}
\fancyhead[R]{\Contributors}

\newcommand{\ProblemDivider}[1]{%
  \vspace{1.0em}
  \noindent\textbf{#1}\par
  \vspace{0.35em}
}


\definecolor{problemBack}{RGB}{235,242,250}
\definecolor{problemFrame}{RGB}{70,110,160}
\definecolor{exerciseGray}{gray}{0.96}
\definecolor{exerciseGrayFrame}{gray}{0.45}
\newtcolorbox{problemBox}{
  colback=exerciseGray,
  colframe=exerciseGrayFrame,
  arc=2mm,
  left=8pt,right=8pt,top=8pt,bottom=8pt
}


\newtcolorbox{solutionBox}{
  colback=white,
  colframe=white,
  arc=0mm,
  left=0pt,right=0pt,top=0pt,bottom=0pt
}

\newcommand{\ProblemAnnounce}{\noindent}
\newcommand{\SolutionAnnounce}{\noindent\textbf{Solution.}\par\medskip}

\newcommand{\Exercise}[1]{\textbf{Exercise~#1.}\;}

\newcommand{\R}{\mathbb{R}}
\newcommand{\C}{\mathbb{C}}
\newcommand{\Pcal}{\mathcal{P}}

\begin{document}

\begin{center}
  {\Large \textbf{\AssignmentType~\AssignmentNumber}}
\end{center}
\vspace{0.75em}

\ProblemDivider{1}

\begin{problemBox}
\ProblemAnnounce
\Exercise{4.10}
For \(p \in \mathcal{P}(\mathbb{R})\) , define \(T_{p} : \mathbb{R} \longrightarrow \mathbb{R}\) by
\[
T_{p}(x) =
\begin{cases}
\dfrac{p(x) - p(3)}{x-3}, & x \neq 3, \\[10pt]
p'(3), & x = 3.
\end{cases}
\]
For each \(x\in \mathbb{R}\):
\begin{enumerate}
  \item[(a)] Suppose \(p(x) = x^{2} - 3x + 2\). Find \(q\in \mathcal{P}(\mathbb{R})\) such that
  \[
  p(x) - p(3) = (x - 3)q(x).
  \]
  \item[(b)] Prove that \(T_{p} = q\).
  \item[(c)] Show that \(T_{p}\in \mathcal{P}(\mathbb{R})\) for every polynomial \(p\in \mathcal{P}(\mathbb{R})\) and also show that \(T\) : \(\mathcal{P}(\mathbb{R})\longrightarrow \mathcal{P}(\mathbb{R})\) is a linear map.
\end{enumerate}
\end{problemBox}

\begin{solutionBox}
\SolutionAnnounce
\doublespacing

\begin{enumerate}
  \item[(a)] Given \(p(x) = x^{2} - 3x + 2\), we compute \(p(3) = 9 - 9 + 2 = 2\). Then
  \[
  p(x) - p(3) = (x^{2} - 3x + 2) - 2 = x^{2} - 3x = x(x - 3).
  \]
  Therefore,
  \[
  p(x) - p(3) = (x - 3) \cdot x,
  \]
  so \(q(x) = x\).

  \item[(b)] For \(x \neq 3\),
  \[
  T_{p}(x) = \frac{p(x) - p(3)}{x-3} = \frac{(x-3)q(x)}{x-3} = q(x) = x.
  \]
  For \(x = 3\),
  \[
  T_{p}(3) = p'(3) = 2 \cdot 3 - 3 = 3 = q(3).
  \]
  Thus \(T_{p}(x) = q(x)\) for all \(x \in \R\), i.e., \(T_{p} = q\).

  \newpage

  \item[(c)] Let \(p \in \Pcal(\R)\). For \(x \neq 3\), \(T_{p}(x) = \frac{p(x) - p(3)}{x-3}\) is a polynomial in \(x\). For \(x = 3\), \(T_{p}(3) = p'(3)\) is a real number, which can be considered a constant polynomial. Thus \(T_{p} \in \Pcal(\R)\) for every \(p \in \Pcal(\R)\).

  Now we show that \(T : \Pcal(\R) \to \Pcal(\R)\) is linear. Let \(f, g \in \Pcal(\R)\) and \(c \in \R\).

  For \(x \neq 3\):
  \begin{align*}
  T(cf)(x) &= \frac{(cf)(x) - (cf)(3)}{x-3} = c \cdot \frac{f(x) - f(3)}{x-3} = c\, Tf(x), \\
  T(f+g)(x) &= \frac{(f+g)(x) - (f+g)(3)}{x-3} = \frac{f(x)-f(3)}{x-3} + \frac{g(x)-g(3)}{x-3} = Tf(x) + Tg(x).
  \end{align*}

  For \(x = 3\):
  \begin{align*}
  T(cf)(3) &= (cf)'(3) = c \cdot f'(3) = c\, Tf(3), \\
  T(f+g)(3) &= (f+g)'(3) = f'(3) + g'(3) = Tf(3) + Tg(3).
  \end{align*}

  Thus \(T\) satisfies both homogeneity and additivity, so \(T\) is a linear map (3.1).
\end{enumerate}

\end{solutionBox}

\ProblemDivider{2}

\begin{problemBox}
\ProblemAnnounce
\Exercise{4.11}
Suppose \(p\in \mathcal{P}(\mathbb{C})\). Define \(q:\mathbb{C}\longrightarrow \mathbb{C}\) by
\[
q(z) = p(z)\overline{p(\overline{z})}.
\]
\begin{enumerate}
  \item[(a)] If \(\alpha ,\beta \in \mathbb{C}\) , show that \(\alpha \overline{\alpha} = |\alpha |^{2}\in \mathbb{R}\) , and \(\alpha \overline{\beta} +\overline{\alpha}\beta \in \mathbb{R}\).
  \item[(b)] If \(p(z) = 2iz - 2 + 3i\) , find \(q(z)\) . Make sure it is a polynomial with real coefficients. [Hint: Set \(\alpha = 2i\) , and \(\beta = -2 + 3i\) . Then rewrite \(p(z)\) as \(\alpha z + \beta\) . Using this expression, find first \(\overline{p(\overline{z})}\) , and then \(q(z)\) .]
  \item[(c)] For arbitrary \(p\in \mathcal{P}(\mathbb{C})\) , prove that \(q\) is a polynomial with real coefficients. [Hint: Prove first that \(\overline{p(\overline{z})}\) is actually a polynomial on \(z\) .]
\end{enumerate}
\end{problemBox}

\begin{solutionBox}
\SolutionAnnounce
\doublespacing

\begin{enumerate}
  \item[(a)] Let \(\alpha = a + bi\) and \(\beta = c + di\) with \(a,b,c,d \in \R\). Then
  \[
  \alpha \overline{\alpha} = (a+bi)(a-bi) = a^{2} + b^{2} = |\alpha|^{2} \in \R.
  \]
  Next,
  \begin{align*}
  \alpha \overline{\beta} + \overline{\alpha} \beta
  &= (a+bi)(c-di) + (a-bi)(c+di) \\
  &= (ac - adi + bci - bdi^{2}) + (ac + adi - bci - bdi^{2}) \\
  &= 2ac + 2bd \in \R.
  \end{align*}

  \item[(b)] Let \(p(z) = 2iz - 2 + 3i = \alpha z + \beta\) with \(\alpha = 2i,\ \beta = -2 + 3i\). Then
  \[
  \overline{p(\overline{z})} = \overline{\alpha \overline{z} + \beta}
  = \overline{\alpha} z + \overline{\beta}
  = (-2i)z + (-2 - 3i).
  \]
  Thus
  \begin{align*}
q(z) &= p(z)\,\overline{p(\overline{z})} \\
     &= (\alpha z + \beta)(\overline{\alpha} z + \overline{\beta}) \\
     &= \alpha\overline{\alpha}\, z^2 + \alpha\overline{\beta}\, z + \overline{\alpha}\beta\, z + \beta\overline{\beta} \\
     &= |\alpha|^2 z^2 + (\alpha\overline{\beta} + \overline{\alpha}\beta)z + |\beta|^2 \\
     &= 4z^2 + \big( (2i)(-2-3i) + (-2i)(-2+3i) \big) z + 13 \\
     &= 4z^2 + 12z + 13.
\end{align*}

  The coefficients \(4, 12, 13\) are all real.

  \item[(c)] Let \(p(z) = a_{0} + a_{1}z + \dots + a_{n}z^{n}\) with \(a_{k} \in \C\). Then
  \[
  \overline{p(\overline{z})}
  = \overline{a_{0} + a_{1}\overline{z} + \dots + a_{n}\overline{z}^{\,n}}
  = \overline{a_{0}} + \overline{a_{1}} z + \dots + \overline{a_{n}} z^{n},
  \]
  which is a polynomial in \(z\) with coefficients \(\overline{a_{k}}\). Thus \(q(z) = p(z)\overline{p(\overline{z})}\) is a polynomial. Its coefficients are sums of terms of the form \(a_{i}\overline{a_{j}}\) or \(a_{i}\overline{a_{j}} + \overline{a_{i}}a_{j}\), which are real by part (a). Therefore, \(q\) is a polynomial with real coefficients.
\end{enumerate}

\end{solutionBox}

\end{document}
